% Auto-generated by experiments/export_latex_tables.py — do not cite numbers manually.
% Source: C:\Users\苏萌\Desktop\FeverFavor\data\cache\eval_results\pipeline_eval.json
\begin{table}[t]
\centering
\caption{端到端整体管线评分（CMCR-Colloquial 分层抽样，LLM 裁判两层量表）。裁判与生成器共用模型池，分数存在自我偏好风险，仅作相对参考。}
\label{tab:pipeline-eval}
\begin{tabular}{lrl}
\hline
\textbf{指标} & \textbf{数值} & \textbf{说明} \\
\hline
样本数（评分成功） & 30 & -- \\
R1-R6 总分均值（100 分制） & 67.9 & -- \\
红旗组总分均值 & 72.0 & 硬指标组 \\
非红旗组总分均值 & 65.9 & -- \\
Veto 1 伪造病例率 (\%) & 0.0 & 须为 0 \\
Veto 2 反向致命结论率 (\%) & 3.3 & 红旗组须为 0 \\
Veto 3 违规处方率 (\%) & 0.0 & 须为 0 \\
Veto 4 捏造患者事实率 (\%) & 10.0 & 防线：双栏拆分+数值守卫 \\
红旗组 Veto 2 计数 & 0 & 硬性指标 = 0 \\
红旗组“立即急诊”判定率 (\%) & 90.0 & 越接近 100 越好 \\
Gold 病例精排 Top-1 命中率 (\%) & 86.7 & 参考指标 \\
引用病例编号有效率 (\%) & 100.0 & Veto 1 确定性校验 \\
LLM 生成模式占比 (\%) & 100.0 & 回落模板的样本计入 template \\
\hline
\multicolumn{3}{l}{\textit{症候群分层均值}} \\
\quad 发热 & 72.2 & n=10 \\
\quad 腹痛 & 59.6 & n=10 \\
\quad 红旗·呼吸困难 & 80.7 & n=3 \\
\quad 红旗·胸痛 & 82.0 & n=3 \\
\quad 红旗·头痛 & 40.0 & n=2 \\
\quad 红旗·意识改变 & 76.0 & n=2 \\
\hline
\end{tabular}
\end{table}
